% Generated draft: requires mathematical and source review.
\documentclass[11pt]{article}
\usepackage[margin=25mm]{geometry}
\usepackage{fontspec}
\setmainfont{Latin Modern Roman}
\usepackage{amsmath,amssymb,mathtools,unicode-math}
\setmathfont{Latin Modern Math}
\usepackage{enumitem,needspace,xurl,hyperref}
\setlength{\parindent}{0pt}
\setlength{\parskip}{5pt}
\setlength{\emergencystretch}{3em}
\newcommand{\N}{\mathbb{N}}
\newcommand{\Z}{\mathbb{Z}}
\newcommand{\Q}{\mathbb{Q}}
\newcommand{\R}{\mathbb{R}}
\newcommand{\C}{\mathbb{C}}
\newcommand{\F}{\mathbb{F}}
\newcommand{\A}{\mathbb{A}}
\newcommand{\PP}{\mathbb{P}}
\newcommand{\E}{\mathbb{E}}
\newcommand{\eps}{\varepsilon}
\newcommand{\dd}{\,\mathrm{d}}
\begin{document}
\section*{UnsolvedMath ID: 6200033}
\textbf{Source problem code:} catalogue entry
\subsection{Definitions and mathematical statement}
Let $G$ be finitely generated and let $\Gamma$ be a Cayley graph for a finite generating set, with its path metric. A metric space has asymptotic dimension at most $m$ if, for every $R>0$, it admits a uniformly bounded cover such that each radius-$R$ ball meets at most $m+1$ cover members. Subsets are given the induced metric.

The conjectural implication in the source is
\[
\operatorname{asdim}G\geq n,\qquad X\subseteq\Gamma,\qquad
\operatorname{asdim}X\leq n-2,\qquad X\text{ coarsely separates }\Gamma
\quad\Longrightarrow\quad
G\text{ splits over }H\leq G\text{ with }\operatorname{asdim}H\leq n-1.
\]
A nontrivial splitting means an amalgamated-product or HNN decomposition. A customary coarse-separation convention requires some closed neighborhood $N_R(X)$ whose complement has at least two components that are deep relative to $X$, meaning neither is contained in any $N_S(X)$. This convention is made explicit for interpretation; the source does not supply its definition. The subgroup dimension may be interpreted using the metric induced from $G$, without assuming $H$ finitely generated.

The natural nonnegative-dimension range is $n\geq2$, but the source gives no range or empty-space convention for smaller $n$. It also does not specify a stronger uniform/deep-component definition of separation or any omitted group hypotheses. No finiteness of $\operatorname{asdim}G$ beyond the written lower bound is added.

The source appends a separate continuum statement: a homogeneous continuum, in its stated sense, is a locally connected compact connected metric space whose homeomorphism group acts transitively; it reports that every simply connected such space has no separating simple arc. A simple arc is an embedded copy of $[0,1]$, and separation means disconnected complement. This is supplied context attributed to Papasoglu, not an extra group hypothesis or a proof of the splitting conjecture.
\subsection{Short English statement}
If a subset two dimensions smaller than a finitely generated group coarsely separates its Cayley graph, must the group split over a subgroup at most one dimension smaller? The source's coarse-separation convention and parameter range need checking.
\subsection{Sources}
\begin{itemize}
\item Local UnsolvedMath registry, permanent ID 6200033, exact source code catalogue entry, ``Boundaries of Groups and Kleinian Groups---Problem 33''; conjecture attributed to Panos Papasoglu, 2005. Registry status open, classification OPEN-TRIAGE.
\item Kapovich, \emph{Problems on Boundaries of Groups and Kleinian Groups} (2005), Problem 33, PDF page 9, \url{https://www.math.ucdavis.edu/~kapovich/EPR/problems.pdf}. Supplied normalized pdf-text extraction accessed July 29, 2026. No PDF was accessed here.
\item The source attributes the appended simply connected homogeneous-continuum no-separating-arc statement to Papasoglu, without a precise reference. The research summary relates the group conjecture to quasi-line coarse-separation splittings and asymptotic-dimension bounds, and reports no general resolution located; it supplies no exact later citation.
\item Provenance flags current status and statement-reuse terms NEEDS\_REVIEW; the document is publicly author-hosted.
\end{itemize}
\end{document}
